\section{Set-series, dreamcatchers, and singularity conservation}
\label{sec:dreamcatchers}

Siegel's 2019 paper encodes invariant sets by holomorphic generating series
and then retains only their dominant radial boundary data.  The resulting
$\ell^2$ object is called a \emph{dreamcatcher}.  This is a different
analytic route from the numen: the numen codes branch compositions, whereas
the permutation operator below acts on coefficient functions and detects
sets invariant under inverse image.

\subsection{Set-series and the historical Hydra class}

For $V\subseteq\N_0$, define
\begin{align}
 \varsigma_V(w)&=\sum_{v\in V}w^v, && |w|<1,\label{eq:ordinary-set-series}\\
 \psi_V(z)&=\varsigma_V(e^{2\pi iz})
            =\sum_{v\in V}e^{2\pi ivz},&&\Im z>0.\label{eq:Fourier-set-series}
\end{align}
The coefficients are the indicator function $\ind{n\in V}$.  The 2019 paper
also uses exponential and Dirichlet set-series.  Its displayed Dirichlet
series $\sum_{v\in V}v^{-s}$ requires either $0\notin V$ or an explicit
treatment of the term $v=0$.

The Hydra used in that paper is the following one-dimensional historical
specialization.  It is not the controlling 2026 definition.

\begin{definition}[2019 $\varrho$-Hydra]
Fix $\varrho\geq2$.  A 2019 $\varrho$-Hydra is a surjection
$H:\N_0\to\N_0$ with
\begin{equation}
\label{eq:2019-hydra}
 H(n)=\frac{a_jn+b_j}{d_j}
 \quad(n\equiv j\pmod\varrho),
\end{equation}
where $a_j,d_j\geq1$, $b_j\geq0$, $\gcd(a_j,d_j)=1$, and each assigned
branch is integer-valued and nonnegative.  Set
\begin{equation}
\label{eq:mu-j}
 \mu_j=\frac{\varrho a_j}{d_j}\in\N.
\end{equation}
It is \emph{prime} when $\varrho$ is prime and every $a_j$ is $1$ or a
prime.  It is \emph{regulated} when $\mu_j=1$ for at least one $j$.
\end{definition}

The integrality of \cref{eq:mu-j} follows by comparing
$H(\varrho m+j)$ with $H(j)$, and coprimality then gives
$d_j\mid\varrho$.  Surjectivity is an extra hypothesis of the 2019 analytic
framework; it is not part of the 2026 word formalism.

\subsection{Permutation operators}

Let $\ell^\infty(\N_0)$ carry its supremum norm.  The coefficient pullback
is the linear map
\[
 \mathsf P_H:\ell^\infty(\N_0)\longrightarrow\ell^\infty(\N_0),
 \qquad (\mathsf P_Hc)_n=c_{H(n)}.
\]
Because the 2019 Hydra $H$ is surjective,
\[
 \|\mathsf P_Hc\|_\infty
 =\sup_{n\geq0}|c_{H(n)}|
 =\sup_{m\geq0}|c_m|=\|c\|_\infty;
\]
thus $\mathsf P_H$ is an isometric injection (not necessarily a surjection).
For $c\in\ell^\infty(\N_0)$, both
\[
 \mathcal G_{\mathbb D}c(w)=\sum_{n\geq0}c_nw^n
 \quad(|w|<1),
 \qquad
 \mathcal G_{\mathbb H}c(z)=\sum_{n\geq0}c_ne^{2\pi inz}
 \quad(\Im z>0)
\]
converge normally on compact subsets of their displayed domains.  The two
analytic operators are therefore the conjugates of this specified
coefficient map on the ranges of the injective generating maps:
\[
 \mathcal Q_H:\operatorname{ran}(\mathcal G_{\mathbb D})
 \longrightarrow\operatorname{ran}(\mathcal G_{\mathbb D}),
 \qquad
 \mathscr Q_H:\operatorname{ran}(\mathcal G_{\mathbb H})
 \longrightarrow\operatorname{ran}(\mathcal G_{\mathbb H}).
\]
\begin{align}
 \mathcal Q_H(\mathcal G_{\mathbb D}c)
   &=\mathcal G_{\mathbb D}(\mathsf P_Hc)
     =\sum_{n\geq0}c_{H(n)}w^n,\label{eq:ordinary-permutation}\\
 \mathscr Q_H(\mathcal G_{\mathbb H}c)
   &=\mathcal G_{\mathbb H}(\mathsf P_Hc)
     =\sum_{n\geq0}c_{H(n)}e^{2\pi inz}.
     \label{eq:Fourier-permutation}
\end{align}
For set-series this simply records inverse image:
\begin{equation}
\label{eq:Q-inverse-image}
 \mathcal Q_H\varsigma_V=\varsigma_{H^{-1}(V)},
 \qquad
 \mathscr Q_H\psi_V=\psi_{H^{-1}(V)}.
\end{equation}
Thus $H^{-1}(V)=V$ if and only if the corresponding set-series is a fixed
point.

Let $\xi_m=e^{2\pi i/m}$.  Splitting the input index modulo $\varrho$ and
using
\[
 H(\varrho n+j)=\mu_jn+H(j)
\]
gives, for every
$\psi\in\operatorname{ran}(\mathcal G_{\mathbb H})$, the finite branch
formula
\begin{equation}
\label{eq:Fourier-permutation-branch}
 (\mathscr Q_H\psi)(z)
 =\sum_{j=0}^{\varrho-1}\frac1{\mu_j}
   \sum_{k=0}^{\mu_j-1}\xi_{\mu_j}^{-kH(j)}
   e^{-2\pi i(b_j/a_j)z}
   \psi\left(\frac{\varrho z+k}{\mu_j}\right).
\end{equation}
To verify the formula, expand the right side.  The inner root-of-unity sum
selects coefficients whose indices are congruent to $H(j)$ modulo $\mu_j$;
the exponential prefactor converts the surviving exponent back to
$\varrho n+j$.

\begin{proposition}[fixed coefficient functions]
\label{prop:permutation-fixed-coefficients}
For a coefficient function $c:\N_0\to\C$, define its coefficient pullback by
\[
 (\mathsf P_Hc)_n=c_{H(n)}.
\]
Then $\mathsf P_Hc=c$ if and only if
\begin{equation}
\label{eq:coeff-orbit-constant}
 c_n=c_{H(n)}\qquad(n\geq0).
\end{equation}
Equivalently, $c$ is constant on every total orbit class of $H$.
\end{proposition}

\begin{proof}
By definition, $\mathsf P_Hc=c$ is exactly
\cref{eq:coeff-orbit-constant}.
Iterating that identity makes $c$ constant along forward orbits; if two
points have a common forward image, their coefficients agree.  Conversely,
constancy on orbit classes implies \cref{eq:coeff-orbit-constant}.
\end{proof}

\begin{sourceissue}[``basis'' in the Permutation Operator Theorem]
The source says that the set-series of all orbit classes form a basis of the
fixed-point space.  \Cref{prop:permutation-fixed-coefficients} gives the exact
coefficient-level statement.  If infinitely many orbit classes have
nonzero coefficients, the resulting decomposition is an infinite sum, not
a finite Hamel-basis expansion.  A topological basis claim requires a named
function-space topology and a convergence proof.
\end{sourceissue}

\subsection{Edgepoints and normalization}

Let $\mathbb H_{+i}=\{z\in\C:\Im z>0\}$ and
$\psi\in\mathcal A(\mathbb H_{+i})$.

\begin{definition}[edgepoint and virtual residue]
An $x\in\R$ is an \emph{edgepoint} of $\psi$ when the finite limit
\[
 \operatorname{VRes}_x[\psi]
 =\lim_{y\downarrow0}y\psi(x+iy)
\]
exists.  It is a \emph{virtual pole} when that limit is nonzero.  Write
\[
 P(\psi)=\{x\in\R:x\text{ is an edgepoint of }\psi\}.
\]
\end{definition}

For a Fourier set-series, the elementary estimate
\[
 |\psi_V(x+iy)|\leq\sum_{n\geq0}e^{-2\pi ny}
 =\frac1{1-e^{-2\pi y}}
\]
implies
\begin{equation}
\label{eq:no-super-simple-growth}
 \lim_{y\downarrow0}y^{1+\delta}\psi_V(x+iy)=0
 \qquad(x\in\R,\ \delta>0).
\end{equation}
Thus the virtual residue records the largest possible radial growth order
for a digital Fourier series.

\begin{sourceissue}[ordinary versus virtual residues]
If $\psi$ has an ordinary simple pole at $x$ with complex residue $A$, then
$y\psi(x+iy)\to A/i=-iA$.  The source's definition of a ``classical virtual
pole'' subtracts $\operatorname{VRes}_x[\psi]/(z-x)$ and thereby omits this
factor of $i$.  This edition retains the term \emph{virtual residue} for the
radial normalization above but does not identify it numerically with the
ordinary complex residue.
\end{sourceissue}

\begin{proposition}[cross-section densities give the virtual residue]
\label{prop:cross-section-vres}
Let $\alpha\geq1$, $\beta\in\Z$, and assume that every limit
\[
 d_k=\lim_{N\to\infty}\frac1N
 |V\cap(k+\alpha\Z)\cap[0,N)|
 \qquad(0\leq k<\alpha)
\]
exists.  Then $\beta/\alpha$ is an edgepoint of $\psi_V$ and
\begin{equation}
\label{eq:VRes-cross-section-density}
 \operatorname{VRes}_{\beta/\alpha}[\psi_V]
 =\frac1{2\pi}\sum_{k=0}^{\alpha-1}
 d_k\xi_\alpha^{\beta k}.
\end{equation}
In particular, if $V$ has natural density $d(V)$, then
$\operatorname{VRes}_0[\psi_V]=d(V)/(2\pi)$.
\end{proposition}

\begin{proof}
Put $a_n=\ind{n\in V}\xi_\alpha^{\beta n}$.  Splitting a Ces\`aro sum by
residue class gives
\[
 \lim_{N\to\infty}\frac1N\sum_{n=0}^{N-1}a_n
 =\sum_{k=0}^{\alpha-1}d_k\xi_\alpha^{\beta k}=:L.
\]
The elementary Abel implication for a bounded sequence now gives
\[
 (1-r)\sum_{n\geq0}a_nr^n\longrightarrow L
 \qquad(r\uparrow1).
\]
For completeness, if $A_N=\sum_{n<N}(a_n-L)$, then $A_N=o(N)$ and summation
by parts gives
$(1-r)\sum_{n\geq0}(a_n-L)r^n=(1-r)^2\sum_{N\geq1}A_Nr^{N-1}\to0$;
split the last sum at a fixed index after bounding $|A_N|/N$ on its tail.
Finally set $r=e^{-2\pi y}$.  Since
$y/(1-e^{-2\pi y})\to1/(2\pi)$,
\[
 y\psi_V(\beta/\alpha+iy)
 =\frac{y}{1-r}(1-r)\sum_{n\geq0}a_nr^n
 \longrightarrow\frac{L}{2\pi},
\]
which proves the assertion.  Only the Abel implication is used; no Tauberian
converse is invoked.
\end{proof}

Dreamcatchers therefore contain residue-class density data, not merely a
list of analytic singularities.

\subsection{The Square Sum Lemma}

Equip subsets of $\R/\Z$ with counting measure.  The result below is stated
in Siegel's paper, which credits the proof idea to accepted MathOverflow
answer 333801 by user \emph{fedja} \cite{FedjaMO333801}.  The argument below
is instead this edition's finite-subset Bessel proof; it is not attributed to
the answerer.  It avoids the source proof's interchange of infinite sums
before summability has been established.

\begin{lemma}[Square Sum Lemma]
\label{lem:square-sum}
Let $V\subseteq\N_0$ and put $P'=P(\psi_V)\cap[0,1)$.  Then
\begin{equation}
\label{eq:square-sum}
 \sum_{t\in P'}|\operatorname{VRes}_t[\psi_V]|^2
 \leq\frac1{(2\pi)^2}.
\end{equation}
In particular, the virtual-residue function belongs to $\ell^2(P')$ and
has at most countable nonzero support.
\end{lemma}

\begin{proof}
Let $a_n=\ind{n\in V}$, let $r=e^{-2\pi y}$, and put on bounded sequences
the semidefinite inner product
\[
 \langle u,v\rangle_y=(1-r)\sum_{n\geq0}u_n\overline{v_n}r^n.
\]
For $t\in\R/\Z$, set $e_t(n)=e^{-2\pi int}$.  Then
\begin{align*}
 \langle e_s,e_t\rangle_y
 &=\frac{1-r}{1-re^{2\pi i(t-s)}}
   \longrightarrow\ind{s=t},\\
 \langle a,e_t\rangle_y
 &=(1-r)\psi_V(t+iy)
   \longrightarrow2\pi\operatorname{VRes}_t[\psi_V],\\
 \|a\|_y^2&=(1-r)\sum_{n\geq0}a_nr^n\leq1.
\end{align*}
Fix a finite $F\subseteq P'$.  The Gram matrix
$G_y=(\langle e_s,e_t\rangle_y)_{s,t\in F}$ tends to the identity and is
invertible for small $y$.  The finite-dimensional projection inequality is
\[
 b_y^*G_y^{-1}b_y\leq\|a\|_y^2,
 \qquad b_y=(\langle a,e_t\rangle_y)_{t\in F}.
\]
Passing to the limit yields
$(2\pi)^2\sum_{t\in F}|\operatorname{VRes}_t[\psi_V]|^2\leq1$.
Taking the supremum over finite $F$ gives \cref{eq:square-sum}.
\end{proof}

\subsection{Dreamcatchers and branch-invariant carriers}

\begin{definition}[dreamcatcher over $T$]
Let $\psi\in\mathcal A(\mathbb H_{+i})$, and let
$T\subseteq P(\psi)$ be $1$-periodic.  The function $\psi$ \emph{admits a
dreamcatcher over $T$} when
\begin{equation}
\label{eq:dreamcatcher}
 R_{\psi,T}(x)=\operatorname{VRes}_x[\psi]
 \quad(x\in T)
\end{equation}
belongs to $\ell^2(T/\Z)$.  Following the source, denote the class of such
functions by $\mathcal D_1(T)$.
\end{definition}

The carrier $T$ is \emph{$H$-branch invariant} when
\begin{equation}
\label{eq:H-branch-invariant}
 \frac{\varrho T+k}{\mu_j}\subseteq T
 \quad(0\leq j<\varrho,\ 0\leq k<\mu_j).
\end{equation}
This ensures that every virtual residue appearing after applying a branch in
\cref{eq:Fourier-permutation-branch} is defined on the same carrier.

\begin{lemma}[Virtual Residue Formula]
\label{lem:VRF}
Let $T$ be $H$-branch invariant and let
$\psi\in\mathcal D_1(T)\cap\operatorname{ran}(\mathcal G_{\mathbb H})$.
At every $t\in T$, the function $\mathscr Q_H\psi$ has an edgepoint and
\begin{equation}
\label{eq:VRF}
 \operatorname{VRes}_t[\mathscr Q_H\psi]
 =\frac1\varrho\sum_{j=0}^{\varrho-1}
 e^{-2\pi i(b_j/a_j)t}
 \sum_{k=0}^{\mu_j-1}\xi_{\mu_j}^{-kH(j)}
 R_{\psi,T}\left(\frac{\varrho t+k}{\mu_j}\right).
\end{equation}
\end{lemma}

\begin{proof}
Multiply \cref{eq:Fourier-permutation-branch} by $y$ and substitute
$z=t+iy$.  The exponential prefactor converges to its value at $t$.  For
$A_{j,k}(z)=(\varrho z+k)/\mu_j$, the imaginary part of $A_{j,k}(t+iy)$ is
$\varrho y/\mu_j$, hence
\[
 y\psi(A_{j,k}(t+iy))
 =\frac{\mu_j}{\varrho}
   \left(\frac{\varrho y}{\mu_j}\right)
   \psi\left(A_{j,k}(t)+i\frac{\varrho y}{\mu_j}\right).
\]
The factor $1/\mu_j$ in \cref{eq:Fourier-permutation-branch} cancels the
$\mu_j$ here, leaving $1/\varrho$.  The sum is finite, so its limit is the
sum of the limits.
\end{proof}

\begin{definition}[dreamcatcher operator]
For $R\in\ell^2(T/\Z)$ let $\mathcal E_{H:T}R:T\to\C$ denote the raw finite
expression
\begin{equation}
\label{eq:dreamcatcher-operator}
 (\mathcal E_{H:T}R)(x)
 =\frac1\varrho\sum_{j=0}^{\varrho-1}
 e^{-2\pi i(b_j/a_j)x}
 \sum_{k=0}^{\mu_j-1}\xi_{\mu_j}^{-kH(j)}
 R\left(\frac{\varrho x+k}{\mu_j}\right).
\end{equation}
Here $R(y)$ means the value on the class $y+\Z$; branch invariance ensures
that every argument belongs to $T$.  Define the maximal graph domain
\[
 \operatorname{Dom}(\mathfrak Q_{H:T})
 =\{R\in\ell^2(T/\Z):
     \mathcal E_{H:T}R\text{ is $1$-periodic on $T$ and }
     [\mathcal E_{H:T}R]\in\ell^2(T/\Z)\}.
\]
The dreamcatcher operator is the linear operator
\[
 \mathfrak Q_{H:T}:\operatorname{Dom}(\mathfrak Q_{H:T})
 \subseteq\ell^2(T/\Z)\longrightarrow\ell^2(T/\Z),
 \qquad R\longmapsto[\mathcal E_{H:T}R].
\]
\end{definition}

\begin{theorem}[Singularity Conservation Law]
\label{thm:SCL}
Let $T$ be $H$-branch invariant, let
$\psi\in\mathcal D_1(T)\cap\operatorname{ran}(\mathcal G_{\mathbb H})$ be
fixed by $\mathscr Q_H$, and put $R=R_{\psi,T}$.  Suppose
$R\in\operatorname{Dom}(\mathfrak Q_{H:T})$.  Then
\begin{equation}
\label{eq:SCL}
 \mathfrak Q_{H:T}R=R.
\end{equation}
\end{theorem}

\begin{proof}
Since $\mathscr Q_H\psi=\psi$, the left side of
\cref{eq:VRF} is $\operatorname{VRes}_t[\psi]=R(t)$.  Its right side is
$\mathcal E_{H:T}R(t)$, whose descended class is
$\mathfrak Q_{H:T}R$.  Equality holds pointwise on $T$ and hence in
$\ell^2(T/\Z)$.
\end{proof}

\begin{sourceissue}[undefined class and operator domain]
The Virtual Residue Formula and SCL in the source assume
$\psi\in\mathcal D_2(T)$, but only $\mathcal D_1(T)$ is defined.  The intended
symbol appears to be $\mathcal D_1(T)$.  More seriously, the proof that
\cref{eq:dreamcatcher-operator} preserves $\ell^2(T/\Z)$ treats the inclusion
in \cref{eq:H-branch-invariant} as saying that each affine branch map is a
bijection of $T$.  A bijection of $\R$ whose image sends $T$ into $T$ need
not map $T$ onto itself.  The typed SCL above therefore uses the exact
maximal graph domain and makes membership an explicit hypothesis for the
particular dreamcatcher.  It does not assert that the graph domain is all of
$\ell^2(T/\Z)$ or that the operator is bounded.  Either conclusion needs
two-sided branch invariance or a separate finite-multiplicity norm estimate.
\end{sourceissue}

A function $\psi$ is \emph{tame} in the source when every rational number is
an edgepoint.  In that case $T=\Q$ is available and, by
\cref{lem:square-sum}, the rational virtual residues of a Fourier set-series
form an element of $\ell^2(\Q/\Z)$.

\subsection{The Finite Dreamcatcher Theorem}

The later half of the 2019 paper studies
$\ker(I-\mathfrak Q_H)$ purely on $\ell^2(\Q/\Z)$.  Its central statement is
the following.

\begin{status}[Finite Dreamcatcher Theorem, source theorem]
\label{status:FDT}
Let $H$ be a prime, regulated 2019 Hydra.  If
\[
 R\in\ker(I-\mathfrak Q_{H:\Q})\subseteq\ell^2(\Q/\Z)
\]
and $\supp R$ is finite modulo $\Z$, then
\begin{equation}
\label{eq:FDT}
 R(x)=R(0)\ind{x\in\Z}
 \qquad(x\in\Q).
\end{equation}
The source proof proceeds through the Off-$\varrho$ Theorem, an auxiliary
functional equation, two support propositions, and the Off-$\varrho$
Infinity Theorem.  This edition records the result with that proof status;
it does not use the earlier flawed $\ell^2(T)$-invariance argument as an
independent proof of the finite-support theorem.
\end{status}

The result says that a finitely supported conserved singularity pattern has
no nonintegral rational frequency.  Combining it with the SCL gives the
paper's first weak rationality consequence.

\begin{status}[Weak Rationality Theorem 1, source consequence]
If $V\subseteq\N_0$ is rational in the sense that its ordinary generating
function is rational, and if $H^{-1}(V)=V$ for a prime regulated Hydra, then
$V$ is finite or cofinite in $\N_0$.
\end{status}

Indeed, rationality makes the virtual-pole support finite modulo $\Z$;
invariance gives a permutation-operator fixed point; SCL and FDT force the
only possible dominant rational boundary frequency to be $0$.  The remaining
step uses the eventual periodicity of a $0$--$1$ sequence with rational
generating function.

\begin{sourceissue}[results not promoted from the 2019 paper]
Three later claims are kept out of the proved chain.
\begin{enumerate}[label=\textup{(\roman*)}]
\item The proof of Weak Rationality Theorem 2 asserts that naturally
      measurable subsets of $\N_0$ form a $\sigma$-algebra and that natural
      density is countably additive.  Both assertions are false in general,
      so that theorem is not canonized as proved.
\item Theorem 12 assumes branch invariance of both the rational edgepoints
      and their rational complement.  The source itself calls this double
      requirement the method's ``Achilles heel''; the edition retains the
      result only with that hypothesis if cited.
\item The assertion that a limit of dreamcatchers is the dreamcatcher of the
      limiting function is explicitly Conjecture 1.  A nearby corollary also
      uses an asymptotic fixed-point premise not present in its statement.
\end{enumerate}
\end{sourceissue}
